\section{Background} \label{sec:background}

  The set of \defn{regular expressions} over $\Sigma$, written
  $\RegExp{\Sigma}$, is the smallest set such that
  $0, 1 \in \RegExp{\Sigma}$, $\Sigma \subseteq \RegExp{\Sigma}$,
  and $(R_1 + R_2), (R_1 \cdot R_2), R^* \in \RegExp{\Sigma}$ for
  all $R, R_1, R_2 \in \RegExp{\Sigma}$.
  $SubExp(R)$ for $R \in \RegExp{\Sigma}$ is the set of
  subexpressions of $R$, including $R$ itself.

\begin{definition}
  % taken more or less verbatim from Algebraic Program Analysis +
  % Fast Algorithms for Solving Path Problems
  The \defn{standard interpretation} is the function
  $\si{\cdot}: \RegExp{\Sigma} \to \textbf{L}$, where $\textbf{L}$
  is the set of regular languages, such that
  $\si{0} = \emptyset$, $\si{1} = \{\epsilon\}$ where $\epsilon$
  is the empty string, $\si{a} = \{a\}$ for all $a \in \Sigma$,
  and
  \begin{align*}
  \si{R_1 + R_2} &\defeq \si{R_1} \cup \si{R_2} \\
  \si{R_1 \cdot R_2} &\defeq \{x_1x_2| x_1 \in \si{R_1}, x_2 \in
                       \si{R_2}\} \\
  \si{R^*} &\defeq \{x_1x_2\dots x_n| x_1,x_2,\dots,x_n \in \si{R}\}
  \end{align*}
  for all $R, R_1, R_2 \in \RegExp{\Sigma}$.
\end{definition}

Let $X$ and $X'$ be sets of variables such that for each
$x\in X$, there is a corresponding $x' \in X'$. As an abuse of
notation, we use $X = X' \defeq \bigwedge_{x \in X} x =
x'$. $\fv{F}$ for a formula $F$ is the set of free variables in
$F$. $\Formula{X, X'}$ denotes the set of formulas such that
$\fv{F} \subseteq X \cup X'$. 

\begin{definition}
  A \defn{control flow graph} $G = \tuple{V, E, r, \mathscr{F}}$
  is a graph with vertices $V$, edges $E$, a root vertex $r$ of
  indegree 0 from which every vertex $v \in V$ is reachable, and
  a function $\mathscr{F}: E \to \Formula{X, X'}$ for some
  $X, X'$. We use $\mathscr{F}_{\tuple{u,v}}(Y, Y')$ to denote
  $\mathscr{F}(\tuple{u,v})[X \mapsto Y, X' \mapsto Y']$. We let
  $\Paths{G}{u}{v}$ be the set of all (possibly non-simple) paths
  from $u$ to $v$ in $G$.
\end{definition}

\begin{definition}
  A \emph{path expression} $P$ of type $\tuple{u, v}$ on
  $G = \tuple{V, E, r, F}$ is a regular expression over
  $E$ such that $\si{P} \subseteq Paths_G(u,v)$ for some
  $u,v \in V$. We denote the set of path expressions on $G$ by
  $PathExp(G)$.
\end{definition}

Suppose $star$ is a function that takes a formula $F$ and the
pair of variable vocabulairies $X$ and $X'$ that $F$ relates and
returns some overapproximation of the reflexive transitive
closure of $F$ over $X$ and $X'$ containing no other free
variables.

\begin{definition}
  The \emph{transition formula interpretation} is the function
  $\tf{\cdot}: PathExp(G) \to \Formula{X, X'}$ where
  $\tf{0} \equiv false$, $\tf{1} \equiv X = X'$,
  $\tf{\tuple{u, v}} \equiv \mathscr{F}_{\tuple{u, v}}(X, X')$,
  and
  \begin{align*}
    \tf{P_a + P_b} &\equiv \tf{P_a} \lor \tf{P_b} \\
    \tf{P_a \cdot P_b} &\equiv \exists X''. \tf{P_a}(X, X'')
                         \land \tf{P_b}(X'', X') \\
    \tf{P^*} &\equiv star(\tf{P}, \tuple{X, X'}) 
  \end{align*}
  for all $P, P_a, P_b \in PathExp(G)$. 

  We use $\tf{P}(Y, Y')$ to denote
  $\tf{P}[X \mapsto Y, X' \mapsto Y']$.
\end{definition}

\subsection{Control flow theory}

\begin{definition}
  In a control flow graph $G = \tuple{V, E, r, \mathscr{F}}$,
  $v\in V$ \emph{dominates} $u \in V$, written $v \dominates u$,
  if every path from $r$ to $u$ in $G$ goes through $v$. If
  $v \dominates u$ and $u \ne v$, then $v$ \emph{strictly
    dominates} $u$, written $v \strictlydominates u$. If $v$
  strictly dominates $u$ and is dominated by every vertex that
  strictly dominates $u$, then $v$ is the \emph{immediate
    dominator} of $u$. Note that every vertex besides $r$ has a
  unique immediate dominator --- we denote this by
  $idom(u)$. $E|_v$ is the set of edges $\tuple{u,w}$ such that
  $v \strictlydominates w$. $\tuple{u, v} \in E$ is a \emph{back
    edge} if $v \dominates u$.
\end{definition}

\begin{definition}
  A control flow graph $G = \tuple{V, E, r, \mathscr{F}}$ is
  \emph{reducible} if for every cycle $v_0,v_1,\dots,v_n=v_0$ in
  $G$, there is some $0 \le i \le n$ such that
  $v_i \dominates v_j$ for all $0 \le j \le n$. (Intuitively,
  every cycle has a natural entrypoint.)
\end{definition}

\textbf{TODO}: missing definitions for
\begin{itemize}
\item eval/link/update data structure
\end{itemize}
